% !TEX root = ../auto-diff/AutoDiff.tex

%%% SECTION PART OF optim-smooth.tex



Gradient evaluation $\nabla f(x)$ is the main operation in batch and stochastic gradient methods. For linear models and shallow networks, explicit formulas usually reduce it to matrix-vector products. For deeper or more elaborate computations, expanded derivative formulas become unwieldy and can duplicate work. Automatic differentiation organizes the chain rule recursively to reuse intermediate results.

%%%%%%%%%%%%%%%%%%%%%%%%%%%%%%%%%%%%%%%%%%%%%%%%%%%%%%%%%%%%%%%%%%%%%%%%
\subsection{Finite Differences and Symbolic Calculus}

For $f : \RR^p \rightarrow \RR$, we want to evaluate $\nabla f : \RR^p \mapsto \RR^p$. A finite-difference approximation with a small parameter $\epsilon>0$ is
\eq{
 	\frac{1}{\epsilon}(f(x+\epsilon \de_1)-f(x), \ldots,f(x+\epsilon \de_p)-f(x))^\top	\approx \nabla f(x)
}
This requires $p+1$ evaluations of $f$. Here $\de_k=(0,\ldots,0,1,0,\ldots,0)$ is the coordinate vector with its $1$ at position $k$.
%
For large $p$, repeated evaluations are expensive. Finite differences also balance truncation error against floating-point cancellation as $\epsilon$ decreases. Reverse-mode automatic differentiation computes the full gradient at a cost typically proportional to one evaluation of $f$.
%
Like symbolic differentiation, automatic differentiation applies exact derivative rules, subject to floating-point error. It organizes those rules around the operations of the program that evaluates the function, preserving reuse of intermediate computations.


%%%%%%%%%%%%%%%%%%%%%%%%%%%%%%%%%%%%%%%%%%%%%%%%%%%%%%%%%%%%%%%%%%%%%%%%
\subsection{Computational Graphs}

Let $x=(x_1,\ldots,x_s)$ be the inputs to a program evaluating $f$. Number its intermediate variables $x_{s+1},\ldots,x_t$, with output $x_t=f(x)$, and let $x_k\in\RR^{n_k}$. We seek the Jacobian blocks $\pd{f(x)}{x_k}\in\RR^{n_t\times n_k}$ for $k=1,\ldots,s$. The scalar case $n_k=1$ is easiest to follow; the formulas also apply to vector variables if the matrix products retain the displayed order. The total input dimension is $p=\sum_{k=1}^s n_k$.

\begin{figure}
\centering
\includegraphics[width=.7\linewidth]{chapter-optimization-and-machine-learning-advanced-topics/compgraph/original/comp-graph.jpg}
\caption{\label{fig-compgraph}
A computational graph.
}
\end{figure}

Represent the computation as a sequence of elementary maps:
\eq{
	\foralls k=s+1,\ldots,t, \quad x_k = f_k( x_1,\ldots,x_{k-1} )
}
Each $f_k$ depends only on earlier variables; see Figure~\ref{fig-compgraph}. The directed acyclic graph (DAG) connects each argument used by $f_k$ to its output $x_k$. The indexing is a topological order: every edge points from a smaller to a larger index.
%
Evaluating $f(x)$ is a forward traversal of the graph.
% 
Automatic differentiation differentiates the supplied computational graph. It does not by itself choose an efficient implementation of the original formula; the derivative computation benefits from the efficiency and reuse already present in the forward program.



%%%%%%%%%%%%%%%%%%%%%%%%%%%%%%%%%%%%%%%%%%%%%%%%%%%%%%%%%%%%%%%%%%%%%%%%
\subsection{Forward Mode of Automatic Differentiation}

Forward mode propagates derivatives of each intermediate variable with respect to a chosen input block, here $x_1$. For scalar inputs, repeating the traversal for all $p$ coordinates gives the full derivative. The formulas below also propagate a block Jacobian when $n_1>1$.

\newcommand{\bk}[1]{\left[#1\right]}
\newcommand{\pdb}[2]{\bk{\pd{#1}{#2}}}

\begin{figure}
\centering
\includegraphics[width=.4\linewidth]{chapter-optimization-and-machine-learning-advanced-topics/forward-backward/original/forward-mode.jpg}
\quad
\includegraphics[width=.4\linewidth]{chapter-optimization-and-machine-learning-advanced-topics/forward-backward/original/backward-mode.jpg}
\caption{\label{fig-forward-backward}
Derivative dependencies in forward mode (left) and reverse mode (right).
}
\end{figure}

Initialize the input derivatives by
\eq{
	\pd{x_1}{x_1} = \Id_{n_1 \times n_1}, \quad 
	\pd{x_2}{x_1} = 0_{n_2 \times n_1}, \ldots, \quad
	\pd{x_s}{x_1} = 0_{n_s \times n_1}, 
}
These are one and zero for scalar inputs. Propagate the derivatives using
\eq{
	\foralls k=s+1,\ldots,t, \quad
	\pd{x_k}{x_1} 
	= \sum_{\ell \in \text{parent}(k)} \bk{\pd{x_k}{x_\ell}} \times \pd{x_\ell}{x_1}
	= \sum_{\ell \in \text{parent}(k)} \pd{f_k}{x_\ell}(x_1,\ldots,x_{k-1}) \times \pd{x_\ell}{x_1}.
}
The notation ``$\text{parent}(k)$'' denotes the nodes $\ell<k$ of the graph that are connected to $k$, see Figure~\ref{fig-forward-backward}, left. 
%
The stored quantities are the derivatives $\pd{x_\ell}{x_1}$. The symbol $\times$ denotes matrix multiplication in the vector-valued case.
%
The brackets $[\ldots]$ indicate a local derivative: $\pd{x_k}{x_\ell}$ is the derivative of the elementary map $f_k$ with respect to that argument. It can be evaluated when needed, assuming the local derivative is available in closed form.

Suppose the local derivative operations $\pd{f_k}{x_\ell}$ have comparable costs, as when the variable dimensions $n_k$ are uniformly bounded, and each node has a bounded number of parents. Computing the full Jacobian coordinate by coordinate requires $p$ traversals, each with cost comparable to a forward evaluation. This becomes expensive for large $p$.



%%%
\paragraph{Simple example.}


We consider the function
\eql{\label{eq-simple-func-autodiff}
 	f(x,y) = y\log(x)+\sqrt{y\log(x)}
} 
on the open domain $x>0$, $y\log x>0$, where all elementary derivatives exist. Its computational graph is displayed in Figure~\ref{fig-dag-example-simple}. The iterations of the forward mode to compute the derivative with respect to $x$ read
\begin{align*}
		\pd{x}{x} &= 1,  \quad \pd{y}{x} = 0 \\
		\pd{a}{x} &= \bk{\pd{a}{x}} \pd{x}{x} = \frac{1}{x} \pd{x}{x} &
			\{x \mapsto a = \log(x)\}\\
		\pd{b}{x} &= \bk{\pd{b}{a}} \pd{a}{x}  + \bk{\pd{b}{y}} \pd{y}{x}  = y \pd{a}{x} + 0 &
			\{(y,a) \mapsto b=ya\}\\
		\pd{c}{x} &= \bk{\pd{c}{b}} \pd{b}{x}  = \frac{1}{2\sqrt{b}} \pd{b}{x}&
			\{b \mapsto c=\sqrt{b}\} \\
		\pd{f}{x} &= \bk{\pd{f}{b}} \pd{b}{x} + \bk{\pd{f}{c}} \pd{c}{x}  = 1\pd{b}{x} + 1\pd{c}{x}&
			\{(b,c) \mapsto f=b+c\}
\end{align*}
A second forward pass gives the derivative with respect to $y$:
\begin{align*}
		\pd{x}{y} &= 0,  \quad \pd{y}{y} = 1 \\
		\pd{a}{y} &= \bk{\pd{a}{x}} \pd{x}{y} = 0 &
			\{x \mapsto a = \log(x)\}\\
		\pd{b}{y} &= \bk{\pd{b}{a}} \pd{a}{y}  + \bk{\pd{b}{y}} \pd{y}{y}  = 0 + a \pd{y}{y}  &
			\{(y,a) \mapsto b=ya\}\\
		\pd{c}{y} &= \bk{\pd{c}{b}} \pd{b}{y}  = \frac{1}{2\sqrt{b}} \pd{b}{y}&
			\{b \mapsto c=\sqrt{b}\} \\
		\pd{f}{y} &= \bk{\pd{f}{b}} \pd{b}{y} + \bk{\pd{f}{c}} \pd{c}{y}  = 1\pd{b}{y} + 1\pd{c}{y}&
			\{(b,c) \mapsto f=b+c\}
\end{align*}



\begin{figure}
\centering
\includegraphics[width=.6\linewidth]{chapter-optimization-and-machine-learning-advanced-topics/dag-example-simple/original/example-simple.png}
\caption{\label{fig-dag-example-simple}
Example of a simple computational graph.
}
\end{figure}

\if 0
%%%
\paragraph{Simple example.}

We consider the function
\eql{\label{eq-simple-func-autodiff}
 	f(x) = \log(x)+\sqrt{\log(x)}
} 
on the open domain $x>0$, $y\log x>0$, where all elementary derivatives exist. Its computational graph is displayed in Figure~\ref{fig-dag-example-simple}. The iterations of the forward mode read
\begin{align*}
		\pd{x}{x} &= 1 \\
		\pd{y}{x} &= \bk{\pd{y}{x}} \pd{x}{x} = \frac{1}{x} \pd{x}{x} &
			\{x \mapsto y = \log(x)\}\\
		\pd{z}{x} &= \bk{\pd{z}{y}} \pd{y}{x} = \frac{1}{2\sqrt{y}} \pd{y}{x} &
			\{y \mapsto z=\sqrt{y}\}\\
		\pd{f}{x} &= \bk{\pd{f}{z}} \pd{z}{x} + \bk{\pd{f}{y}} \pd{y}{x} = 1 \pd{z}{x} + 1 \pd{y}{x}&
			\{(x,z) \mapsto f=x+z\}
\end{align*}



\begin{figure}
\centering
\includegraphics[width=.6\linewidth]{auto-diff/example-complex}
\caption{\label{fig-dag-example-complex}
Example of a more complex computational graph.
}
\end{figure}

%%%
\paragraph{More complex example.}

We now consider the function
\eql{\label{eq-cpx-func-autodiff}
	f(x) = \frac{\log(x+\sqrt{x^2+1})}{x^2} - \frac{\log^3(x+\sqrt{x^2+1})}{\sqrt{x^2+1}}
}
whose computational graph is displayed on Figure~\ref{fig-dag-example-complex}. The iterations of the forward mode read
\begin{align*}
		\pd{x}{x} &= 1 \\
		\pd{y}{x} &= \pdb{y}{x} \pd{x}{x}  = 2 x \pd{x}{x}  &
			\{x \mapsto y = x^2\} \\
		\pd{z}{x} &= \pdb{z}{y} \pd{y}{x} = \frac{1}{2\sqrt{y+1}} \pd{y}{x} &
			\{y \mapsto z=\sqrt{y+1}\}\\
		\pd{r}{x} &= \pdb{r}{z} \pd{x}{x} + \pdb{r}{z} \pd{z}{x} = \frac{1}{x+r} \pd{x}{x} + \frac{1}{x+r} \pd{z}{x}&
			\{(x,z) \mapsto r=\log(x+z)\} \\
		\pd{s}{x} &= \pdb{s}{r} \pd{r}{x} + \pdb{s}{y}\pd{y}{x} = \frac{1}{y} \pd{r}{x} - \frac{r}{y^2}\pd{y}{x} &
			\{ (r,y) \mapsto s=\frac{r}{y}  \} \\
		\pd{t}{x} &= \pdb{t}{r} \pd{r}{x} = 3r^2 \pd{r}{x} &
			\{ r \mapsto t = r^3 \} \\
		\pd{u}{x} &= \pdb{u}{t} \pd{t}{x} + \pdb{u}{z}\pd{z}{x} = \frac{1}{z} \pd{t}{x} - \frac{t}{z^2} \pd{z}{x} &
			\{ (t,z) \mapsto u =  \frac{t}{z}  \} \\
		\pd{f}{x} &= \pdb{f}{s}\pd{s}{x} + \pdb{f}{u}\pd{u}{x} = 1 \pd{s}{x} - 1 \pd{u}{x} &
			\{ (s,u) \mapsto f = s-u \}
\end{align*}
\fi


%%%
\paragraph{Dual numbers.}

Dual numbers provide a convenient implementation of the forward pass. They form a real algebra whose elements have the form $x + \epsilon x'$ where $\epsilon$ is a symbol obeying the rule that $\epsilon^2=0$. Here $(x,x') \in \RR^2$ and $x'$ is intended to store a derivative with respect to some input variable. Multiplication and, when $x\neq0$, inversion are given by
\eq{
	(x+\epsilon x')(y+\epsilon y') = xy + \epsilon (x y' + y x')
	\qandq
	\frac{1}{x+\epsilon x'} = \frac{1}{x}-\epsilon \frac{x'}{x^2}. 
}
For polynomial or rational $f$, the arithmetic rules imply
\eq{
	f(x+\epsilon) = f(x) + \epsilon f'(x).
}
For another differentiable elementary function $f$, define its action on dual numbers by
\eq{
	f(x+\epsilon x') \eqdef f(x) + \epsilon f'(x) x'.
}
This definition gives
\eq{
	(f \circ g)(x+\epsilon ) = f(g(x)) + \epsilon f'(g(x))g'(x)
}
which is the chain rule. More generally, for $f(x_1,\ldots,x_s)$ built from these elementary operations,
\eq{
	f(x_1+\epsilon,x_2,\ldots,x_s) = f(x_1,\ldots,x_s) + \epsilon \pd{f}{x_1}(x_1,\ldots,x_s)
}
the dual-number evaluation computes $\pd{f}{x_1}(x_1,\ldots,x_s)$ by forward-mode differentiation. The other coordinates are treated identically.


%%%%%%%%%%%%%%%%%%%%%%%%%%%%%%%%%%%%%%%%%%%%%%%%%%%%%%%%%%%%%%%%%%%%%%%%
\subsection{Reverse Mode of Automatic Differentiation}\label{sec-reverse-mode}

Reverse mode computes the derivatives $\pd{x_t}{x_k}$ of the output with respect to the intermediate variables. For a scalar output and many input coordinates, this avoids the repeated forward traversals required to compute each $\pd{x_k}{x_i}$ separately.



Initialize the derivative of the output node by
\eq{
	\pd{x_t}{x_t} = \Id_{n_t \times n_t}, 
}
and apply the following recursion from the last node to the first:
\eq{
	\foralls k=t-1,t-2,\ldots,1, \quad
	\pd{x_t}{x_k} 
	= \sum_{m \in \text{son}(k)}  \pd{x_t}{x_m} \times \bk{\pd{x_m}{x_k}}
	= \sum_{m \in \text{son}(k)} \pd{x_t}{x_m} \times \pd{f_m(x_1,\ldots,x_{m-1})}{x_k}.
}
%
The set $\text{son}(k)$ contains the children $m>k$ that depend directly on $x_k$, as shown in Figure~\ref{fig-forward-backward}, right. Contributions from every child must be added.

%%%
\paragraph{Back-propagation.}

For a scalar output $x_t \in \RR$, we have $\pd{x_t}{x_k} = [\nabla_{x_k} f(x)]^\top \in \RR^{1 \times n_k}$. The reverse recursion can then be written for gradient vectors:
\eq{
	\foralls k=t-1,t-2,\ldots,1, \quad
	\nabla_{x_k} f(x) 
	= \sum_{m \in \text{son}(k)} \pa{\pd{f_m(x_1,\ldots,x_{m-1})}{x_k}}^\top \pa{
	 	\nabla_{x_m} f(x)
	}.
}
Here $\pa{\pd{f_m(x_1,\ldots,x_{m-1})}{x_k}}^\top \in \RR^{n_k \times n_m}$ is the adjoint Jacobian of $f_m$. This adjoint recursion is called backpropagation and is the standard form used in machine learning.

For a scalar output ($n_t=1$), reverse mode is usually the preferred way to compute the full gradient. It requires access to the intermediate variables $(x_k)_{k=s+1}^t$, which can be costly to store for a large graph. Checkpointing trades additional computation for reduced storage.

%%%
\paragraph{Simple example.}

For the two-input example $f(x,y)$ in~\eqref{eq-simple-func-autodiff}, reverse mode gives
\begin{align*}
		\pd{f}{f} &= 1 &\\
		\pd{f}{c} &= \pd{f}{f} \pdb{f}{c} = \pd{f}{f} 1 &
			\{ c \mapsto f = b+c\}\\
		\pd{f}{b} &= \pd{f}{c} \pdb{c}{b} + \pd{f}{f} \pdb{f}{b} = \pd{f}{c} \frac{1}{2\sqrt{b}} + \pd{f}{f} 1 & 
			\{ b \mapsto c=\sqrt{b}, b \mapsto f=b+c\} \\
		\pd{f}{a} &= \pd{f}{b} \pdb{b}{a} = \pd{f}{b} y & 
			\{ a \mapsto b=y a\} \\
		\pd{f}{y} &= \pd{f}{b} \pdb{b}{y} = \pd{f}{b} a & 
			\{ y \mapsto b=ya\} \\
		\pd{f}{x} &= \pd{f}{a} \pdb{a}{x} = \pd{f}{a} \frac{1}{x} & 
			\{ x \mapsto a=\log(x)\}
\end{align*}
A single reverse traversal computes both derivatives with respect to $x$ and $y$, whereas coordinatewise forward mode requires two passes.

\if 0
%%%
\paragraph{More complex example.}

We consider once again the function $f(x)$ of~\eqref{eq-cpx-func-autodiff}, the iterations of the reverse mode read
\begin{align*}
		\pd{f}{f} &= 1 &\\
		\pd{f}{u} &= \pd{f}{f} \pdb{f}{u} 
			= \pd{f}{f} (-1) &
			\{u \mapsto f = s-u\}\\
		\pd{f}{t} &= \pd{f}{u} \pdb{u}{t} 
			= \pd{f}{u} \frac{1}{z}  & 
			\{ t \mapsto u= \sqrt{y}, y \frac{t}{z} \} \\
		\pd{f}{s} &= \pd{f}{f} \pdb{f}{s} = \pd{f}{f} 1 & 
			\{ s \mapsto f = s-u \} \\
		\pd{f}{r} &= \pd{f}{t} \pdb{t}{r} + \pd{f}{s} \pdb{s}{r} 
			= \pd{f}{t} 3r^2 + \pd{f}{s} \frac{1}{y} & 
			\{ r \mapsto s = \frac{r}{y}, r \mapsto  t=r^3\}  \\
		\pd{f}{z} &= \pd{f}{u} \pdb{u}{z} + \pd{f}{r} \pdb{r}{z} 
			= \pd{f}{u} \frac{-t}{z^2} + \pd{f}{r} \pdb{1}{x+z} &
			\{ z \mapsto u=\frac{t}{z}, z \mapsto r=\log(x+z)\} \\
		\pd{f}{y} &= \pd{f}{z}\pdb{z}{y} + \pd{f}{s}\pdb{s}{y}
			= \pd{f}{z} \frac{1}{2\sqrt{y+1}} + \pd{f}{s}\frac{-r}{y^2} &
			\{ y \mapsto z=\sqrt{y+1}, y \mapsto s = \frac{r}{y}\} \\
		\pd{f}{x} &= \pd{f}{y}\pdb{y}{x} + \pd{f}{r}\pdb{r}{x}
			= \pd{f}{y} 2x + \pd{f}{r} \frac{1}{x+z}  &
			\{x \mapsto y=x^2, x \mapsto r=\log(x+z)\}		
\end{align*}
\fi


%%%%%%%%%%%%%%%%%%%%%%%%%%%%%%%%%%%%%%%%%%%%%%%%%%%%%%%%%%%%%%%%%%%%%%%%
\subsection{Feed-forward Compositions}\label{sec-reversemode-simple}

For a feedforward chain, the computed function is the composition
\eql{\label{eq-simple-lin-dag}
	f = f_{t} \circ f_{t-1} \circ \ldots \circ f_2 \circ f_1
}
for functions $f_k : \RR^{n_{k-1}} \rightarrow \RR^{n_{k}}$.

The forward evaluation starts from $x_0=x \in \RR^{n_0}$ and computes
\eq{
	\foralls k=1, \ldots, t, \quad x_{k} = f_k(x_{k-1})
}
giving the output $f(x) = x_t$.

Let $A_k \eqdef \partial f_k(x_{k-1}) \in \RR^{n_{k} \times n_{k-1}}$ be the Jacobian of the local map. The chain rule gives
\eq{
	\partial f(x) = A_{t} \times A_{t-1} \times \ldots A_2 \times A_1.
}
Forward mode multiplies the Jacobians from right to left, while reverse mode multiplies them from left to right
\begin{align*}
	\partial f(x) &= A_{t} \times \pa{  A_{t-1} \times \pa{ \ldots \times \pa{ A_3 \times \pa{ A_2 \times A_1 } } } }, \\
	\partial f(x) &= \pa{ \pa{ \pa{ \pa{ A_{t} \times A_{t-1} }  \times A_{t-2} } \times \ldots } \times A_2 } \times A_1.
\end{align*}


\begin{figure}
\centering
\includegraphics[width=.42\linewidth]{chapter-optimization-and-machine-learning-advanced-topics/matrix-mult/original/matrix-forward.jpg}
\qquad
\includegraphics[width=.45\linewidth]{chapter-optimization-and-machine-learning-advanced-topics/matrix-mult/original/matrix-backward.jpg}
\caption{\label{fig-matrix-mult}
Matrix multiplication costs for forward mode (left) and reverse mode (right) along a chain.
}
\end{figure}

We note that the computation of the product $A \times B$ of $A \in \RR^{n \times p}$ with $B \in \RR^{p \times q}$ requires $O(npq)$ arithmetic operations with the standard dense algorithm.
%
As shown on Figure~\ref{fig-matrix-mult}, the arithmetic costs of the two parenthesizations are proportional to
\eq{
	n_0 \sum_{k=1}^{t-1} n_k n_{k+1}
	\qandq
	n_t \sum_{k=0}^{t-2} n_k n_{k+1}
}
For roughly comparable intermediate widths, these costs favor reverse mode when $n_t\ll n_0$. A scalar learning objective typically has $n_t=1$, making reverse mode a natural choice; the exact cost comparison depends on all the intermediate dimensions.


%%%%%%%%%%%%%%%%%%%%%%%%%%%%%%%%%%%%%%%%%%%%%%%%%%%%%%%%%%%%%%%%%%%%%%%%
\subsection{Feed-forward Architecture}

A parameterized feedforward architecture, such as a neural network, has the form
\eql{\label{eq-feednets}
	\foralls k=1, \ldots, t, \quad x_{k} = f_k(x_{k-1},\th_{k-1})
}
where $\th_{k-1}$ is the parameter vector for a layer and the input $x_0 \in \RR^{n_0}$ is fixed. The objective is
\eql{\label{eq-loss-feedf}
	f(\th) \eqdef L(x_t)
}
Here $L : \RR^{n_t} \rightarrow \RR$ is a loss, such as squared prediction error or logistic loss, and $\th=(\th_k)_{k=0}^{t-1}$ collects the parameters. Figure~\ref{fig-mlp} shows the computational graph.


\begin{figure}
\centering
\includegraphics[width=.7\linewidth]{chapter-optimization-and-machine-learning-advanced-topics/mlp/original/archi-feedforward.jpg}
\caption{\label{fig-mlp}
Computational graph for a feedforward architecture.
}
\end{figure}


Reverse mode computes the gradient of $f$ by propagating derivatives through the variables $(x_k,\th_k)$. Initialize
\eq{
	\nabla_{x_t} f = \nabla L(x_t) 	
}
and recurse from $k=t$ down to $1$:
\eql{\label{eq-backprop-discr}
	z_{k-1} = [\partial_x f_k(x_{k-1},\th_{k-1})]^\top z_k
	\qandq
	\nabla_{\th_{k-1}} f = [\partial_\th f_k(x_{k-1},\th_{k-1})]^\top (\nabla_{x_{k}} f)
}
where $z_k \eqdef \nabla_{x_{k}} f(\th)$ is the gradient with respect to $x_{k}$.

%%%
\paragraph{Multilayer perceptrons.}\label{sec-autodiff-mlp}

For a fully connected network without biases, use
\eql{\label{eq-mlp-func}
	\foralls x_{k-1} \in \RR^{n_{k-1}}, \quad
	f_k(x_{k-1},\th_{k-1}) = \rho( \th_{k-1} x_{k-1} )
} 
where $\th_{k-1}\in\RR^{n_k\times n_{k-1}}$ contains the weights and $\rho$ acts coordinatewise; see Figure~\ref{fig-mlp-param}. The derivative formulas below apply where $\rho$ is differentiable. At a ReLU kink, a chosen derivative convention defines the computed update, but the classical derivative need not exist.
%
For an adjoint vector $z_k \in \RR^{n_k}$, typically $\nabla_{x_{k}} f$,
\eq{
	\left\{
	\begin{matrix}
	[\partial_x f_k(x_{k-1},\th_{k-1})]^\top (z_k) = \th_{k-1}^\top w_k z_k, \\
	[\partial_\th f_k(x_{k-1},\th_{k-1})]^\top (z_k) =(w_kz_k)x_{k-1}^\top
	\end{matrix}
	\right.
	\qwhereq
	w_k \eqdef \diag(\rho'(\th_{k-1} x_{k-1})).
}

\begin{figure}
\centering
\includegraphics[width=.6\linewidth]{chapter-optimization-and-machine-learning-advanced-topics/mlp-param/original/network-mlp.jpg}
\caption{\label{fig-mlp-param}
Multilayer-perceptron parameterization.
}
\end{figure}



%%%
\paragraph{Link with adjoint state method.}

A residual form of~\eqref{eq-feednets} can be interpreted as a time discretization of an ordinary differential equation.
%
Keep the state dimension fixed, $n_k=n$, and interpret $x_k$ as an approximation of a continuous trajectory $x(t)$ at time $t=k\tau$. The formal continuum limit takes $\tau\downarrow0$ with $k\tau$ approaching a fixed time.
%
Impose the residual structure
\eql{\label{eq-flow-eq}
	f_k(x_{k-1},\th_{k-1}) = x_{k-1} + \tau u(x_{k-1},\th_{k-1}, k\tau)
}
with a parameterized vector field $u(x,\th,t) \in \RR^n$. Formally, as $\tau \rightarrow 0$, this approaches the nonlinear ODE
\eql{\label{eq-ode-structure}
	\dot x(t) = u(x(t),\th(t),t)
}
with $x(t=0)=x_0$.

Writing $z(t) = \nabla_{x(t)} f(\th)$ for the adjoint, the discrete relations~\eqref{eq-backprop-discr} become a linear backward ODE:
\eq{
	\dot z(t)  =  - [\partial_x u(x(t),\th(t),t)]^\top z(t)
	\qandq
	\nabla_{\th(t)} f(\th) = [\partial_\th u(x(t),\th(t),t)]^\top z(t).
}
The adjoint has terminal condition $z(T)=\nabla L(x(T))$. The parameter derivative is interpreted as the $L^2$ functional gradient with respect to the time-dependent control. Its discrete normalization is $\frac{1}{\tau} \nabla_{\th_{k-1}} f  \rightarrow \nabla_{\th(t)} f(\th)$.

%%%%%%%%%%%%%%%%%%%%%%%%%%%%%%%%%%%%%%%%%%%%%%%%%%%
\subsection{Recurrent Architectures}

A recurrent architecture reuses the parameter $\th=\th_k$ and map $f_k=h$ at every step of~\eqref{eq-feednets}:
\eql{\label{eq-feednets-recur}
	\foralls k=1, \ldots, t, \quad x_{k} = h(x_{k-1},\th).
}
Consider a scalar objective
\eq{
	f(\th) = L(x_t,\th)
}
The loss now depends explicitly on $\th$, extending~\eqref{eq-loss-feedf}. Figure~\ref{fig-recur} shows the corresponding graph.


\begin{figure}
\centering
\includegraphics[width=.6\linewidth]{chapter-optimization-and-machine-learning-advanced-topics/recur/original/archi-recurrent.jpg}
\caption{\label{fig-recur}
Computational graph for a recurrent architecture.
}
\end{figure}

Backpropagation gives
\eql{\label{eq-backprop-recurrent}
	\nabla_{x_{k-1}} f = [\partial_x h(x_{k-1},\th)]^\top \nabla_{x_{k}} f
	\qandq
	\nabla_{\th} f = 
		\nabla_\th L(x_t,\th) + 
		\sum_k [\partial_\th h(x_{k-1},\th)]^\top \nabla_{x_{k}} f.
}
Similarly, with $h(x,\th) = x + \tau u(x,\th)$, the formal limit $(k,k\tau) \rightarrow (+\infty,t)$ gives a forward nonlinear ODE with a time-independent vector field:
\eq{
	\dot x(t) = u(x(t),\th)
}
and the following linear backward adjoint equation, for $f(\th)=L(x(T),\th)$
\eql{\label{eq-jacobian-mlp}
	\dot z(t)  =  - [\partial_x u(x(t),\th)]^\top z(t)
	\qandq
	\nabla_{\th} f(\th) = \nabla_\th L(x(T),\th) + \int_0^T[\partial_\th u(x(t),\th)]^\top z(t) \d t.
}
with terminal condition $z(T)=\nabla_xL(x(T),\th)$.

%%%%%
\paragraph{Residual recurrent networks. }

A residual recurrent network uses the update
\eq{
	h(x,\th) = x + W_2^\top \rho( W_1 x )
}
as illustrated in Figure~\ref{fig-recurrent-param}. Here $\th = (W_1, W_2) \in (\RR^{q\times n})^2$ contains the weights and $\rho$ is a pointwise activation. Increasing the hidden width $q$ enlarges the class of representable residual maps.
%
When $W_2=-\tau W_1$ and $\rho=\psi'$, this update is a gradient step for $\Ee(x,\th)=\sum_i\psi((W_1x)_i)$. Repeating it gives an iterative implementation of the argmin layer~\eqref{eq-argmin-layer}, provided the iterations converge to a minimizer.
%
The Jacobians $\partial_\th h$ and $\partial_x h$ follow from the layer formulas in the \hyperref[sec-autodiff-mlp]{multilayer-perceptron discussion}.

\begin{figure}
\centering
\includegraphics[width=.5\linewidth]{chapter-optimization-and-machine-learning-advanced-topics/recurrent-param/original/network-recur.jpg}
\caption{\label{fig-recurrent-param}
Recurrent residual perceptron parameterization.
}
\end{figure}


%%%%%
\paragraph{Mitigating Memory Requirements. }

Backpropagation requires access to the iterates $(x_k)_{k=0}^t$, whose storage can dominate memory use. Checkpointing stores selected states and reruns parts of the forward computation to reconstruct missing values. Recursive schedules can use $O(\log t)$ stored states with an $O(\log t)$ computational overhead.

If the forward computation is invertible, stored states can sometimes be reconstructed instead. Assume maps $g_k$ satisfy
\eq{
	x_{k} = g_k(x_{k+1},\ldots,x_t).
}
These inverse updates may amplify numerical errors, so algebraic invertibility alone does not ensure stable reconstruction. In practice, the inverse can also depend on a few extra variables, in particular on the input values $(x_0,\ldots,x_s)$.

For example, split the continuous state as $x(t)=(r(t),s(t))$ and suppose the vector field in~\eqref{eq-ode-structure} has the separated form $u((r,s),\th,t)=(F(s,\th,t),G(r,\th,t))$. An invertible staggered update is
\eq{
	r_{k+1} = r_k + \tau F(s_k,\th_k,\tau k)
	\qandq
	s_{k+1} = s_k + \tau G(r_{k+1},\th_{k+1/2},\tau (k+1/2)).
}
In exact arithmetic, these updates can be inverted as
\eq{
	s_k = s_{k+1} - \tau G(r_{k+1},\th_{k+1/2},\tau (k+1/2)).
	\qandq
	r_k = r_{k+1}  - \tau F(s_k,\th_k,\tau k). 
}

%%%%%
\paragraph{Fixed point maps}

If $h$ is continuous and the iterates converge to $x^\star(\th)$, their limit is a fixed point:
\eq{
	x^\star(\th) = h(x^\star(\th),\th). 
}
Backpropagation differentiates the finite computation $f(\th) = L(x_t,\th)$. For the limiting objective $f^\star(\th) = L(x^\star(\th),\th)$, the implicit function theorem gives an alternative:
\eql{\label{eq-impl-func-formula}
	\nabla f^\star(\th) = [\partial x^\star(\th)]^\top( \nabla_x L(x^\star(\th),\th) )
		 + \nabla_\th L(x^\star(\th),\th).
}
Assume $h$ is continuously differentiable and $\Id-\partial_xh$ is invertible at the fixed point. The implicit function theorem then gives
\eq{
	\partial x^\star(\th)=\left(\Id-\partial_xh(x^\star(\th),\th)\right)^{-1}\partial_\th h(x^\star(\th),\th).
}
In practice, one replaces in these formulas $x^\star(\th)$ by $x_t$, which approximates the derivative of the converged objective $f^\star$, rather than the derivative of the finite unrolling $f$.
%
This approach replaces storage of the iteration history by a linear solve at the approximate equilibrium.



%%%%%
\paragraph{Argmin layers}

An argmin layer maps a parameter $\th$ to a minimizer $x(\th)$ of a parameterized objective. Assume a unique minimizer, or specify a branch when several exist:
\eq{
	x(\th) = \uargmin{x} \Ee(x,\th).
}
One way to compute the layer is gradient descent, initialized here at $x_0=0$:
\eql{\label{eq-argmin-layer}
	x_{k+1} = x_k - \tau \nabla \Ee(x_k,\th). 
}
This has the form~\eqref{eq-flow-eq} when using the vector field $u(x,\th)=-\nabla_x\Ee(x,\th)$.

Assume $\Ee$ is twice continuously differentiable near the selected minimizer and its Hessian with respect to $x$ is invertible. The implicit function theorem gives a locally differentiable minimizer branch; applying~\eqref{eq-impl-func-formula} to its stationarity equation yields
\eq{
	\nabla f^\star(\th) = -\pa{ \frac{\partial^2 \Ee}{\partial x \partial \th}(x^\star(\th),\th) }^\top 
		\pa{  \pdd{\Ee}{x}(x^\star(\th),\th) }^{-1}
		( \nabla_x L(x^\star(\th),\th) )
		 + \nabla_\th L(x^\star(\th),\th)
}

If $f(\th)$ is the minimized value itself, $f(\th)=\Ee(x^\star(\th),\th)$, then $L=\Ee$. Formula~\eqref{eq-impl-func-formula} simplifies because $\nabla_x L(x^\star(\th),\th)=0$, giving
\eql{\label{eq-danskin}
	\nabla f^\star(\th) = \nabla_\th L(x^\star(\th),\th). 
}
This is the envelope theorem, also associated with Danskin's theorem.


%%%%%
\paragraph{Sinkhorn's algorithm}

Let $a\in\RR^n$ and $b\in\RR^m$ be strictly positive probability vectors, let $C\in\RR^{n\times m}$ be a cost matrix, and let $\epsilon>0$. Entropic optimal transport minimizes
\eq{\sum_{ij}C_{ij}P_{ij}+\epsilon\sum_{ij}P_{ij}(\log P_{ij}-1)}
over $P\geq0$ with row sums $a$ and column sums $b$. Set $K_{ij}=e^{-C_{ij}/\epsilon}>0$. Sinkhorn's alternating updates, initialized with positive scaling vectors, are
\eq{u_{k+1}=\frac{a}{Kv_k},\qquad v_{k+1}=\frac{b}{K^\top u_{k+1}}.}
All divisions are componentwise. The second update uses the newly computed $u_{k+1}$; this sequential dependence is part of Sinkhorn's alternating scheme.

The regularized transport value is the maximum of the dual objective
\eq{\Ee(u,v;a,b)=\epsilon\left(\dotp a{\log u}+\dotp b{\log v}-\dotp u{Kv}\right).}
Sinkhorn performs alternating maximization of this objective. A finite number of iterations defines a differentiable recurrent map on positive inputs. Its Jacobian products can be computed by chaining the two scaling steps. For example, for $u=a/(Kv)$,
\eq{\partial_vu=-\diag\left(\frac{a}{(Kv)^2}\right)K,\qquad
\partial_au=\diag\left(\frac1{Kv}\right),}
and analogous formulas apply to the $v$ update, with the dependence through the newly computed $u$ included.

At convergence, the envelope theorem gives marginal derivatives represented by the dual potentials
\eq{\nabla_a f^\star=\epsilon\log u^\star,\qquad
\nabla_b f^\star=\epsilon\log v^\star.}
The scalings have the invariance $(u,v)\mapsto(cu,v/c)$, so these potentials are defined up to opposite additive constants. On the probability simplices, derivatives act on zero-sum perturbations and are therefore independent of this gauge. An implicit differentiation formula must fix the gauge before inverting a Jacobian; the unmodified fixed-point Jacobian is singular in that direction.
